\chapter{Hardware and Software Tools Used}
\label{ch:tools}

The implementation of the hybrid spatial-frequency image classification system requires a high-performance computing environment capable of handling large datasets and complex spectral computations. This chapter outlines the hardware and software stack used for the development, training, and evaluation of the project.

\section{Hardware Requirements}
The training of deep spatial-spectral models, particularly when processing high-resolution facial images and large-scale object benchmarks, is computationally intensive. The following hardware configuration was utilized to ensure efficient feature extraction and model optimization:

\begin{itemize}
    \item \textbf{GPU:} NVIDIA Tesla V100 (32GB VRAM) – utilized for hardware-accelerated tensor computations in the PyTorch branch and parallelizing the training of the ResNet18 backbone.
    \item \textbf{CPU:} Intel(R) Xeon(R) Gold Multi-core Processor (minimum 12 cores) – essential for the high-speed extraction of 2D-FFT features and handling multi-threaded data loading.
    \item \textbf{RAM:} 64 GB – required to manage high-resolution image buffers, spectral feature matrices, and PCA-based dimensionality reduction operations.
    \item \textbf{Storage:} High-speed NVMe SSD for fast access to raw binary datasets (STL-10) and the storage of model artifacts and logs.
\end{itemize}

\section{Software Requirements}
The software stack includes a combination of deep learning frameworks, classical machine learning libraries, and specialized image processing tools.

\subsection{Programming Language}
\begin{itemize}
    \item \textbf{Python 3.8+:} Chosen as the primary language due to its industry-standard status in AI research and its rich ecosystem of libraries for spectral analysis and machine learning.
\end{itemize}

\subsection{Libraries and Frameworks}
\begin{itemize}
    \item \textbf{PyTorch:} Used for implementing the deep spatial branch (ResNet18) and managing the GPU-accelerated training pipeline.
    \item \textbf{XGBoost:} Utilized for the spectral branch to provide high-performance gradient boosting on PCA-reduced frequency features.
    \item \textbf{NumPy \& SciPy:} Essential for the mathematical implementation of the Fast Fourier Transform (FFT) and handling multidimensional arrays.
    \item \textbf{Scikit-Learn:} Used for Principal Component Analysis (PCA), dataset splitting, and performance evaluation metrics.
    \item \textbf{OpenCV \& Scikit-Image:} Employed for image preprocessing tasks, including CLAHE (Histogram Equalization), Gaussian blurring, and normalization.
    \item \textbf{Grad-CAM (pytorch-grad-cam):} Used for generating visual explainability heatmaps to validate model focus.
    \item \textbf{PyYAML:} For managing centralized configuration files (`.yaml`) to ensure reproducible experiments.
\end{itemize}

\subsection{Development and Deployment Environment}
\begin{itemize}
    \item \textbf{Operating System:} Linux (Ubuntu 20.04/22.04 LTS) – the standard platform for NVIDIA driver support and CUDA-enabled workflows.
    \item \textbf{CUDA Toolkit:} Version 11.x/12.x to enable low-level GPU acceleration for PyTorch and XGBoost.
    \item \textbf{Remote Access Tools:} WinSCP and PuTTY (or SSH) were used for migrating code and datasets to the remote V100 server and managing background training sessions.
\end{itemize}

\subsection{Supporting Tools}
\begin{itemize}
    \item \textbf{Git:} For version control and tracking changes throughout the development lifecycle.
    \item \textbf{Conda:} Used to create isolated virtual environments, ensuring that library dependencies did not conflict.
\end{itemize}

\section{Why These Tools Were Chosen}
The selection of this specific stack was driven by the hybrid nature of the project:

\begin{enumerate}
    \item \textbf{NVIDIA V100 GPU:} The 32GB of VRAM was critical for training ResNet18 and custom FFTNets on large batch sizes without hitting memory limits.
    \item \textbf{XGBoost over standard Neural Networks:} XGBoost was selected for the spectral branch because it is often more effective at handling structured, PCA-reduced data than deep networks, providing a robust second opinion in the hybrid architecture.
    \item \textbf{PyTorch:} Selected for its dynamic nature and excellent integration with CUDA, making it easy to implement custom layers and explainability hooks.
    \item \textbf{WinSCP/SSH:} Essential for a professional workflow, allowing for local development in an IDE (like VS Code) while leveraging the power of a remote server for the heavy computational work.
\end{enumerate}
